\section{Introduction}

Predictive systems often have more auxiliary information available than they
should use at once. A traffic forecaster can consult neighbouring sensors, a
language model can retrieve passages, an in-context classifier can choose
demonstrations, and a tabular model can acquire additional feature groups. In
each case, selecting a useful candidate becomes harder after other candidates
have already been chosen, because the predictor has changed.

We model this changing value through \emph{state-conditioned marginal
utility}: the expected reduction in task loss produced by adding one candidate
to the current context set. The definition shifts the focus from candidate
content to the prediction that already exists. A highly relevant candidate can
be redundant after similar information has been selected. A less obvious
candidate can become valuable when it corrects what the current prediction is
still missing.

The predictor itself exposes a clue about this value. Let
$d_{A,j}=f_{A\cup\{j\}}-f_A$ denote the change in prediction caused by
candidate $j$ at selected set $A$. We call this change the predictor response.
Our theory shows that marginal utility is governed by this response across a
broad class of losses. Smooth losses admit a response expansion with a
controlled remainder, Bregman losses give an exact identity, softmax
cross-entropy has an explicit curvature bound, and squared loss yields an affine
response model once the current state is fixed. These results also reveal the
quantity a selector must learn: the systematic residual left by the fixed
predictor.

This structure leads to a simple routing design. A cached model screens the
candidate pool, the fixed predictor is evaluated only on a small shortlist, and
a response-aware scorer estimates which retained candidate will reduce the
current loss most. We also study a theory-shaped bilinear head that represents
the response--residual interaction directly. The same analysis suggests
measurements for deciding whether routing is worth learning: how much signal is
present in the cached screen, whether candidate order can be recovered from
decision-time observations, and how much information the candidate pool adds
beyond the anchor.

The experiments tell a consistent story across prediction regimes. With a
linear subset predictor, response-aware routing is the strongest deployable
selector in detailed comparisons with content, fixed-set, utility-based and
learned alternatives. With stronger subset-capable predictors, the oracle still
finds useful context combinations, yet the candidate ordering becomes difficult
to recover from the available observations. Increasing the candidate library
widens this gap: the oracle opportunity grows while the learned router loses
rankability. A separate pool-geometry study shows why relevance can switch from
harmful to useful as candidate novelty rises, and the resulting mechanism gives
a label-free adaptive fallback. Demonstration selection on image benchmarks
provides a second task family and separates the value of ranking from the
calibration of stopping.

The paper contributes a unified decision object, a general response--utility
theory, response-aware and theory-shaped selectors, and an empirical account of
what makes their gains attainable. The central message is practical. Context
selection works best when the predictor leaves a systematic, observable
residual that candidates can correct. Predictor strength, pool geometry and
selection rules all change that signal, so a deployment should choose its
selection strategy with those three factors in view.
